\documentclass[11pt]{article}
\usepackage[T1]{fontenc}
\usepackage[utf8]{inputenc}
\usepackage{lmodern}
\usepackage[margin=1in]{geometry}
\usepackage{amsmath,amssymb,mathtools}
\usepackage{cancel}
\usepackage{microtype}
\usepackage[hidelinks]{hyperref}
\setlength{\parindent}{0pt}
\setlength{\parskip}{0.5em}
\allowdisplaybreaks
\newcommand{\C}{\mathbb C}
\newcommand{\R}{\mathbb R}
\newcommand{\Q}{\mathbb Q}
\newcommand{\Z}{\mathbb Z}
\newcommand{\Hd}{\mathcal H_d}
\newcommand{\Prim}{\mathcal P}
\newcommand{\dR}{\mathrm{dR}}
\newcommand{\id}{\mathrm{id}}
\newcommand{\bdel}{\bar\partial}
\newcommand{\orthsum}{\mathbin{\oplus^{\perp}}}
\newcommand{\Vol}{\operatorname{Vol}}
\title{I. Hodge Theory and Periods}
\author{}
\date{}
\begin{document}
\maketitle

% Slide 1.
$(X,\omega)$ compact K\"ahler, $\dim_{\C}X=n$, $L=[\omega]\smile(-)$.

\paragraph{Hodge decomposition.}
\[
 H^k(X;\C)=\bigoplus_{p+q=k}H^{p,q}(X).
\]
Cohomology splits by complex type.

\paragraph{Lefschetz decomposition and Hard Lefschetz.}
\[
 H^k=\bigoplus_{r\geq0}L^rP^{k-2r},
 \qquad L^{n-k}:H^k\xrightarrow{\sim}H^{2n-k}
 \quad(0\leq k\leq n).
\]
\paragraph{Elliptic PDE.}
Unique harmonic representatives $+$ K\"ahler identities.

\paragraph{Hodge--Riemann relations.}
\[
 i^{p-q}Q_{\omega,k}(\alpha,\bar\alpha)>0,
 \qquad 0\neq\alpha\in P^{p,q},\quad p+q=k\leq n.
\]
\paragraph{Polarized pure Hodge structures.}
$(P^k(X;\R),Q_{\omega,k})$ of weight $k$. Over $\Q$ if $[\omega]$ is rational.

% Slide 2.
\section*{Setup and complex types}
$X$: a complex manifold, with its complex orientation.
\[
 J:TX\longrightarrow TX,\qquad J^2=-\id.
\]
The complex structure. $g$: a Riemannian metric.
$A^k=A^k(X;\C)$: smooth complex-valued $k$-forms.

Complex charts determine $J$:
\[
 \begin{pmatrix} I&0\\0&-I\end{pmatrix},
 \qquad T_{\C}X=T^{1,0}X\oplus T^{0,1}X,
\]
\[
 T^{1,0}X=\ker(J-i\id),\qquad
 T^{0,1}X=\ker(J+i\id).
\]
Dualize and take exterior powers:
\[
 A^k=\bigoplus_{p+q=k}A^{p,q}.
\]
Locally, $d(dz_j)=d(d\bar z_j)=0$. Differentiate a $(p,q)$-form
$\Rightarrow$ a $(p+1,q)$- or $(p,q+1)$-form:
\[
 \partial\alpha:=(d\alpha)^{p+1,q},\qquad
 \bdel\alpha=(d\alpha)^{p,q+1}.
\]
\[
 d=\partial+\bdel,\qquad
 \partial^2=\bdel^2=0,\qquad
 \partial\bdel+\bdel\partial=0.
\]

% Slides 3--4.
\section*{Hermitian metrics and the fundamental form}
\paragraph{Hermitian invariance.}
A Riemannian metric $g$ is \textbf{Hermitian} if
\[
 g(Ju,Jv)=g(u,v)\quad\text{for all }u,v\in T_xX.
\]
Consequently, by $J^2$, $g(Ju,v)=-g(u,Jv)$. Put
\[
 \omega(u,v):=g(Ju,v).
\]
The fundamental form $\omega$ is a real $2$-form, of type $(1,1)$,
and positive: $\omega(u,Ju)=g(u,u)>0$ for $u\neq0$.
\[
 \omega(v,u)=g(Jv,u)=g(u,Jv)=-g(Ju,v)=-\omega(u,v).
\]
\[
 \{\text{Hermitian metrics}\}
 \xleftrightarrow{\;1:1\;}
 \{\text{positive real $(1,1)$-forms}\},
\]
\[
 g\longmapsto\bigl(\omega(u,v)=g(Ju,v)\bigr),\qquad
 g(u,v)=\omega(u,Jv).
\]

\subsection*{Hermitian metrics in holomorphic coordinates}
In holomorphic coordinates $z_1,\ldots,z_n$,
\[
 g=2\operatorname{Re}\left(\sum_{j,k=1}^n
 h_{j\bar k}\,dz_j\otimes d\bar z_k\right).
\]
The matrix $(h_{j\bar k})$ is Hermitian and positive definite.
The fundamental form is
\[
 \omega=i\sum_{j,k=1}^n h_{j\bar k}\,dz_j\wedge d\bar z_k.
\]
Positive definiteness permits diagonalization and normalization.
Hence a \textbf{unitary basis} of $T_x^{1,0}X$ exists at every $x\in X$.

% Slide 5.
\section*{The K\"ahler condition and the volume form}
The Hermitian manifold $(X,J,g)$ is \textbf{K\"ahler} if $d\omega=0$.
Then $[\omega]\in H^2_{\dR}(X;\R)$ is the K\"ahler class.
\[
 dV_g=\frac{\omega^n}{n!},\qquad n=\dim_{\C}X.
\]
If $X$ is compact,
\[
 0<\Vol_g(X)=\int_X\frac{\omega^n}{n!}<\infty.
\]

% Slides 6--7.
\section*{Examples of K\"ahler manifolds}
\paragraph{1. Complex curves.}
Every Hermitian metric on a complex curve is K\"ahler:
$d\omega$ would be a $3$-form on a real surface.

\paragraph{2. Flat complex tori.}
The Euclidean metric and form on $\C^n$ are invariant under translations.
They descend to every complex torus
\[
 X=\C^n/\Gamma,\qquad \Gamma\subset\C^n,
\]
producing a flat K\"ahler metric.
\[
 \omega_0=\sum_{j=1}^n dx_j\wedge dy_j,\qquad
 g_0=\sum_{j=1}^n(dx_j^2+dy_j^2).
\]

\paragraph{3. Projective space.}
\[
 \mathbf P^n=(\C^{n+1}\setminus\{0\})/\C^{\times},
 \qquad [Z_0:\cdots:Z_n].
\]
On the chart $Z_0\neq0$, put $z^j=Z_j/Z_0$.
\[
 \omega_{\mathrm{FS}}=i\partial\bdel
 \log\left(1+\sum_{j=1}^n|z^j|^2\right).
\]
Thus these local forms define the Fubini--Study K\"ahler form on
$\mathbf P^n$. Its restriction to a complex submanifold is again closed
and positive.

\paragraph{Complex submanifolds.}
For a complex submanifold $i:Y\hookrightarrow X$, restrict the K\"ahler
form of $X$: $\omega_Y:=i^*\omega_X$.

\paragraph{Products.}
If $(X_1,\omega_1)$ and $(X_2,\omega_2)$ are K\"ahler, their product is
K\"ahler with
\[
 \omega=\pi_1^*\omega_1+\pi_2^*\omega_2,
 \qquad\pi_a:X_1\times X_2\longrightarrow X_a.
\]

% Slide 8. Handwritten expressions are retained as written.
\subsection*{A non-example: the Hopf surface}
\[
 X:=(\C^2\setminus\{0\})/\langle z\mapsto2z\rangle.
\]
The free, properly discontinuous action is holomorphic, so $X$ is a
complex surface.
\[
 X\cong_{\mathrm{diff}}S^1\times S^3.
\]
Write $z=ru$, $r>0$, $u\in S^3$:
\[
 z\longmapsto(\log r,u)\in\R^3\times S^3,
 \qquad z\longmapsto2z\ \leadsto\ \text{translation}.
\]
\[
 H^2(X;\R)=0,\qquad \omega=d\eta,\qquad
 \frac{\omega^2}{2}=dV_g\geq0.
\]
\[
 d\omega=0,\qquad
 \omega^2=d\eta\wedge\omega=d(\eta\wedge\omega).
\]
By Stokes,
\[
 \int_X\omega^2=\int_Xd(\eta\wedge\omega)=0.
\]
\[
 \int_X\omega^n>0,\qquad [\omega]\neq0,\qquad H^2(X;\R)=0
 \quad\Longrightarrow\quad X\text{ not K\"ahler}.
\]

% Slide 9.
\section*{Towards the analytic Hodge theorem: the Hodge star}
$X$: real dimension $2n$.
\[
 *:A^k(X;\R)\longrightarrow A^{2n-k}(X;\R).
\]
For a fixed $k$-form $\beta$, the form $*\beta$ is the unique
$(2n-k)$-form such that
\[
 \alpha\wedge *\beta=\langle\alpha,\beta\rangle_g\,dV_g
 \qquad\text{for all }\alpha\in A^k(X;\R).
\]
\textbf{Exercise.} $X$ compact K\"ahler surface: $*\omega=\omega$.

% Slide 10.
\subsection*{The inner product, adjoints, and Laplacians}
Assume henceforth that $X$ is compact and without boundary.
The Hermitian $L^2$ inner product is
\[
 \langle\alpha,\beta\rangle_{L^2}=\int_X\alpha\wedge *\bar\beta,
\]
linear in $\alpha$. For $D\in\{d,\partial,\bdel\}$, its formal adjoint
$D^*$ is characterized by
\[
 \langle D\alpha,\beta\rangle_{L^2}
 =\langle\alpha,D^*\beta\rangle_{L^2}.
\]
For $D\in\{d,\partial,\bdel\}$, put $\Delta_D=DD^*+D^*D$:
\[
 \Delta_d=dd^*+d^*d,\qquad
 \Delta_{\partial}=\partial\partial^*+\partial^*\partial,\qquad
 \Delta_{\bdel}=\bdel\bdel^*+\bdel^*\bdel.
\]
A form $\alpha\in A^k$ is \textbf{$D$-harmonic} if $\Delta_D\alpha=0$.
Set
\[
 \mathcal H_D^k:=\ker(\Delta_D:A^k\longrightarrow A^k),
\]
the space of $D$-harmonic forms.

% Slide 11.
\subsection*{Harmonic forms are closed}
\paragraph{The energy identity.}
For $D\in\{d,\partial,\bdel\}$,
\[
 \Delta_D=D^*D+DD^*,\qquad
 \langle\Delta_D\alpha,\alpha\rangle_{L^2}
 =\|D\alpha\|_{L^2}^2+\|D^*\alpha\|_{L^2}^2.
\]
\[
 \alpha\text{ is $D$-harmonic}
 \quad\Longleftrightarrow\quad
 D\alpha=0,\quad D^*\alpha=0.
\]
\paragraph{Review.}
$X$ compact, complex, without boundary.
$J\leadsto A^k=\bigoplus_{p+q=k}A^{p,q}$.
The Hermitian metric $(X,J,g)$ gives the Hodge star $*$ and
$\langle\cdot,\cdot\rangle_{L^2}$, then adjoints $D^*$ and Laplacians
$\Delta_D$. $D$-harmonic forms are $D$-closed.
Additional condition: K\"ahler, $d\omega=0$.

% Slide 12.
\subsection*{The analytic Hodge theorem}
On a compact oriented Riemannian manifold without boundary,
\[
 A^k=\Hd^k\orthsum dA^{k-1}\orthsum d^*A^{k+1},
\]
with orthogonality in $L^2$.
The harmonic space is finite-dimensional, and every de Rham class has a
\textbf{unique harmonic representative}.
\[
 \phi_d^k:\Hd^k\xrightarrow{\sim}H^k_{\dR}(X;\C),
 \qquad\alpha\longmapsto[\alpha].
\]
\textbf{Black box:} elliptic operator theory.
\textbf{Goal:} split the harmonic space into Hodge types.

% Slides 13--14.
\section*{The K\"ahler Laplacian preserves bidegree}
Now assume $X$ is compact K\"ahler. For $k=p+q$, define
\[
 \Hd^{p,q}:=\Hd^k\cap A^{p,q}.
\]
% The following definition is crossed out on slide 13.
\[
 \cancel{\mathcal H_{\bdel}^{p,q}
 :=\ker(\Delta_{\bdel}:A^{p,q}\longrightarrow A^{p,q})}.
\]
Different types are orthogonal: $A^{p,q}\perp A^{r,s}$ if
$(p,q)\neq(r,s)$.
Restrict the Hodge map:
\[
 \phi_d^{p,q}:\Hd^{p,q}\xrightarrow{\sim}H^{p,q}_{\dR}(X).
\]
Here $H^{p,q}_{\dR}(X)$ consists of classes admitting a $d$-closed
$(p,q)$-representative.

\textbf{Injectivity:} the analytic Hodge theorem.
\textbf{Surjectivity:} $\Delta_{\bdel}$ preserves type, so the K\"ahler
identity implies that $\Delta_d$ preserves type. Therefore the harmonic
projection $P:A^k\longrightarrow\Hd^k$ also preserves type.

\subsection*{Surjectivity: the harmonic representative has the same type}
Take $[\alpha]\in H^{p,q}_{\dR}(X)$ with $\alpha\in A^{p,q}$ and
$d\alpha=0$. Put $k=p+q$.
Write $Z_d^{p,q}:=\ker(d:A^{p,q}\longrightarrow A^{k+1})$.
The representative and its class are related by
\[
 Z_d^{p,q}\xrightarrow{P}\Hd^{p,q}
 \xrightarrow{\phi_d^{p,q}}H^{p,q}_{\dR}(X),\qquad
 \phi_d^{p,q}\circ P=[\,\cdot\,].
\]
The analytic Hodge theorem gives a unique $h\in\Hd^k$ with
$[h]=[\alpha]$. Since $h=P\alpha$ and $P$ preserves type,
\[
 h\in\Hd^{p,q},\qquad [\alpha]=\phi_d^{p,q}(h).
\]

\paragraph{Why harmonic projection preserves type.}
Let $\pi_{p,q}$ be projection onto the $(p,q)$-part.
\[
 \pi_{p,q}:A^k\longrightarrow A^{p,q},\qquad
 \pi_{p,q}:\Hd^k\longrightarrow\Hd^{p,q},\qquad
 P\pi_{p,q}=\pi_{p,q}P.
\]
$\Delta_{\bdel}$ preserves type and $\Delta_d=2\Delta_{\bdel}$.
Thus $\Delta_d$ preserves type. Its kernel splits into the mutually
orthogonal type subspaces, so orthogonal projection onto the kernel
commutes with $\pi_{p,q}$.
\[
 A^k=\Hd^k\orthsum(\Hd^k)^\perp\qquad(L^2).
\]

% Slide 15.
\subsection*{Splitting the harmonic space}
Let $\alpha\in\Hd^k$ and write
$\alpha=\sum_{p+q=k}\alpha^{p,q}$. Since $\Delta_d$ preserves types,
\[
 0=\Delta_d\alpha=\sum_{p+q=k}\Delta_d\alpha^{p,q}
 \quad\Longrightarrow\quad
 \Delta_d\alpha^{p,q}=0\quad\text{for all }p,q.
\]
Hence
\[
 \Hd^k=\bigoplus_{p+q=k}\Hd^{p,q}.
\]
The Hodge decomposition follows through the isomorphisms
\begin{align*}
 H^k_{\mathrm{sing}}(X;\C)
 &\xrightarrow[\text{de Rham}]{\sim}H^k_{\dR}(X;\C)\\
 &\xrightarrow[\text{analytic Hodge}]{\sim}\Hd^k(X)\\
 &\xrightarrow[\Delta_d\text{ preserves type}]{\sim}
   \bigoplus_{p+q=k}\Hd^{p,q}(X)\\
 &\xrightarrow[\text{restricted Hodge}]{\sim}
   \bigoplus_{p+q=k}H^{p,q}_{\dR}(X).
\end{align*}

% Slide 16.
\section*{Hodge diamond and symmetry}
\[
 h^{p,q}:=\dim_{\C}H^{p,q}_{\dR}(X).
\]
Schematic for $n=3$:
\[
\begin{array}{ccccccc}
 &&&h^{0,0}&&&\\
 &&h^{0,1}&&h^{1,0}&&\\
 &h^{0,2}&&h^{1,1}&&h^{2,0}&\\
 h^{0,3}&&h^{1,2}&&h^{2,1}&&h^{3,0}\\
 &h^{1,3}&&h^{2,2}&&h^{3,1}&\\
 &&h^{2,3}&&h^{3,2}&&\\
 &&&h^{3,3}&&&
\end{array}
\]
Left--right: complex conjugation. Top--bottom: Hodge star.

\paragraph{Complex conjugation.}
$\Delta_d$ is real, hence commutes with conjugation:
\[
 \overline{\Hd^{p,q}}=\Hd^{q,p},\qquad h^{p,q}=h^{q,p}.
\]
\paragraph{Hodge star and analytic Hodge theory.}
The adjoint formulas give $[*,\Delta_d]=0$, so
\[
 *:\Hd^{p,q}\xrightarrow{\sim}\Hd^{n-q,n-p},\qquad
 h^{p,q}=h^{n-q,n-p}.
\]
Together: $h^{p,q}=h^{n-p,n-q}$.

% Slide 17.
\subsection*{Examples and geometric consequences}
For a compact K\"ahler manifold, Hodge decomposition gives
\[
 b_k(X)=\sum_{p+q=k}h^{p,q}(X),\qquad b_{2m+1}(X)\in2\Z.
\]
Here $b_k(X)=\dim_{\R}H^k(X;\R)$. In odd degree, $p\neq q$,
so conjugation pairs all the summands.

\paragraph{The Hopf surface.}
\[
 S:=(\C^2\setminus\{0\})/\langle z\mapsto2z\rangle
 \cong_{\mathrm{diff}}S^3\times S^1.
\]
By the K\"unneth formula,
\[
 H^1(S;\R)\cong H^1(S^1;\R)\cong\R,\qquad b_1(S)=1.
\]
A compact K\"ahler manifold has even $b_1$.
Thus the Hopf surface $S$ is \textbf{not K\"ahler}.

% Slide 18.
\section*{K\"ahler identities}
On a compact K\"ahler $n$-fold, put
\[
 L\alpha=\omega\wedge\alpha,\qquad \Lambda=L^*,\qquad
 [A,B]=AB-BA.
\]
\paragraph{Commutators and their adjoints.}
\[
 [\Lambda,\partial]=i\bdel^*,\qquad
 [\Lambda,\bdel]=-i\partial^*,
\]
\[
 [\partial^*,L]=-i\bdel,\qquad
 [\bdel^*,L]=i\partial.
\]
Since $d\omega=0$, and by adjointness,
\[
 [L,\partial]=[L,\bdel]=0,\qquad
 [\Lambda,\partial^*]=[\Lambda,\bdel^*]=0.
\]
\paragraph{Laplacian consequences.}
\[
 \partial\bdel^*+\bdel^*\partial=0,\qquad
 \bdel\partial^*+\partial^*\bdel=0,
\]
\[
 \Delta_d=2\Delta_{\partial}=2\Delta_{\bdel},\qquad
 [L,\Delta_d]=[\Lambda,\Delta_d]=0.
\]
Thus $L$ and $\Lambda$ preserve harmonic forms.
\paragraph{Pointwise Lefschetz algebra.}
\[
 [\Lambda,L]|_{A^k}=(n-k)\id.
\]

% Slide 19.
\section*{Hard Lefschetz}
\[
 L:A^k\longrightarrow A^{k+2},\qquad\alpha\longmapsto\omega\wedge\alpha.
\]
$\omega$ is closed and has even degree:
\[
 d(L\alpha)=d\omega\wedge\alpha+\omega\wedge d\alpha=L(d\alpha).
\]
Thus $L$ preserves closed forms.
\[
 L(\alpha+d\beta)-L(\alpha)=d(\omega\wedge\beta).
\]
Consequently,
\[
 L:H^k_{\dR}(X;\C)\longrightarrow H^{k+2}_{\dR}(X;\C),
 \qquad[\alpha]_{\dR}\longmapsto[\omega\wedge\alpha]_{\dR}.
\]
In the Hodge diamond for $n=3$, $L$ moves vertically.
Reflection is realized by powers of $L$:
\[
 L^{n-p-q}:H^{p,q}_{\dR}\xrightarrow{\sim}H^{n-q,n-p}_{\dR},
 \qquad p+q\leq n.
\]
$L$ is injective toward the middle ($p+q<n$).
\textbf{Goal:}
\[
 L^{n-k}:H^k(X;\C)\xrightarrow{\sim}H^{2n-k}(X;\C),
 \qquad0\leq k\leq n.
\]

% Slide 20.
\subsection*{Hard Lefschetz: proof strategy}
Study harmonic forms and transfer to cohomology by analytic Hodge:
\[
 \phi_d^k:\Hd^k\xrightarrow{\sim}H^k_{\dR}(X;\C).
\]
Primitive decompositions and the primitive-string calculation:
\[
 \Hd^k\cong\bigoplus_{r=0}^{\lfloor k/2\rfloor}L^r\Prim^{k-2r},
 \qquad
 \Hd^{2n-k}\cong
 \bigoplus_{r=0}^{\lfloor k/2\rfloor}L^{n-k+r}\Prim^{k-2r},
\]
\[
 L^{n-k}:\Hd^k\xrightarrow{\sim}\Hd^{2n-k}
 \qquad(0\leq k\leq n).
\]
What is $\Prim^k$? $\Hd^k=\Prim^k\orthsum L\Hd^{k-2}$.
Orthogonality in $L^2$; iterate.

% Slide 21.
\subsection*{Ingredients for primitive forms and the orthogonal splitting}
$\Lambda:=L^*$ with respect to the $L^2$ Hermitian inner product.
The K\"ahler identities
\[
 [\Lambda,\partial]=i\bdel^*,\qquad
 [\Lambda,\bdel]=-i\partial^*,\qquad
 \Delta_d=2\Delta_{\bdel},
\]
together with $d\omega=0$ and the type $(1,1)$ of $\omega$, imply
\[
 [L,\Delta_{\bdel}]=[L,\Delta_d]=[\Lambda,\Delta_d]=0.
\]
Operators on harmonic forms:
\[
 L:\Hd^k\longrightarrow\Hd^{k+2},\quad\text{bidegree }(1,1),
\]
\[
 \Lambda:\Hd^k\longrightarrow\Hd^{k-2},\quad\text{bidegree }(-1,-1).
\]
$L$ and $\Lambda$ remain adjoint with respect to $L^2$ on the harmonic spaces.

% Slide 22.
\subsection*{Primitive harmonic forms: the first splitting}
\[
 \Prim^k:=\ker(\Lambda:\Hd^k\longrightarrow\Hd^{k-2}).
\]
An element $\alpha\in\Prim^k$ is a \textbf{primitive harmonic form}.
On the finite-dimensional harmonic spaces, $\Lambda=L^*$, so
\[
 \Prim^k=(L\Hd^{k-2})^\perp.
\]
Consequently,
\[
 \Hd^k=\Prim^k\orthsum L\Hd^{k-2}
 \quad\text{with respect to }L^2.
\]
Next: iterate this splitting.

% Slide 23.
\subsection*{Primitive harmonic forms: lengths of the strings}
\[
 \Prim^j=\ker(\Lambda:\Hd^j\longrightarrow\Hd^{j-2}).
\]
Let $0\neq\alpha\in\Prim^j$, with $0\leq j\leq n$.
The degree-$j$ form $\alpha$ generates the nonzero string
\[
 \alpha,L\alpha,\ldots,L^{n-j}\alpha.
\]
\textbf{Claim:} $L^{n-j+1}\alpha=0$.
\[
 \|L^r\alpha\|_{L^2}^2=\boxed{\,?\,}\,\|\alpha\|_{L^2}^2>0,
 \qquad
 L^r:\Prim^j\hookrightarrow\Hd^{j+2r}
 \quad(0\leq r\leq n-j).
\]
\paragraph{Proof.}
On $j$-forms, $[\Lambda,L]=(n-j)\id$.
\[
 \Lambda\alpha=0
 \quad\Longrightarrow\quad
 (L\Lambda+(n-j)\id)\alpha=(n-j)\alpha.
\]
For every $r\geq1$,
\[
 \Lambda L^r\alpha=r(n-j-r+1)L^{r-1}\alpha.
\]
Adjointness gives
\[
 \|L^r\alpha\|_{L^2}^2
 =r(n-j-r+1)\|L^{r-1}\alpha\|_{L^2}^2.
\]
For $r=1$, $\|L\alpha\|_{L^2}^2=(n-j)\|\alpha\|_{L^2}^2$.
So $\Prim^j=0$ for all $j>n$.

% Slide 24.
\subsection*{Primitive strings and the Lefschetz decomposition}
For $\alpha\in\Prim^j$, the case $r=1$ gives
\[
 \|L\alpha\|_{L^2}^2=(n-j)\|\alpha\|_{L^2}^2.
\]
Hence $\Prim^j=0$ for $j>n$.
For $0\neq\alpha\in\Prim^j$ and $0\leq r\leq n-j$,
\[
 \|L^r\alpha\|_{L^2}^2
 =\frac{r!\,(n-j)!}{(n-j-r)!}\|\alpha\|_{L^2}^2>0.
\]
\[
 L^r:\Prim^j\longrightarrow\Hd^{j+2r}
 \quad\text{is injective for }0\leq r\leq n-j.
\]
Iterating the splitting, for a fixed $0\leq j\leq n$,
\[
 \Hd^j=\mathop{\bigoplus\!{}^{\perp}}_{r=0}^{\lfloor j/2\rfloor}
 L^r\Prim^{j-2r}\qquad\text{(orthogonal in }L^2\text{)}.
\]
For primitive $a,b$ in the summands and $r<s$, adjointness and the
lowering formula give
\[
 \langle L^ra,L^sb\rangle
 =\langle a,\Lambda^rL^sb\rangle
 =c\langle a,L^{s-r}b\rangle=0.
\]

% Slide 25.
\subsection*{A picture of the Lefschetz strings}
Primitive families, illustrated for $n=4$.
Each arrow is $L$ and raises cohomological degree by $2$.
\[
\renewcommand{\arraystretch}{1.35}
\begin{array}{c|ccccccccc}
\text{degree}&0&1&2&3&4=n&5&6&7&8\\\hline
 &\Prim^0&&L\Prim^0&&L^2\Prim^0&&L^3\Prim^0&&L^4\Prim^0\\
 &&\Prim^1&&L\Prim^1&&L^2\Prim^1&&L^3\Prim^1&\\
 &&&\Prim^2&&L\Prim^2&&L^2\Prim^2&&\\
 &&&&\Prim^3&&L\Prim^3&&&\\
 &&&&&\Prim^4&&&&
\end{array}
\]
\[
 \Hd^4=L^2\Prim^0\orthsum L\Prim^2\orthsum\Prim^4
 \qquad(L^2).
\]
For $0\neq\alpha\in\Prim^j$, the string
$\alpha,L\alpha,\ldots,L^{n-j}\alpha$ has $n-j+1$ nonzero terms.
Its endpoint degrees are $j$ and $2n-j$; every string is centered at
degree $n$. Each node is a subspace (possibly zero); each row is a
family of primitive strings.

% Slide 26. The handwritten H^2(X;C)=0 is preserved as written.
\subsection*{A smooth quintic threefold}
Let $X=\{F_5=0\}\subset\mathbf P^4$ be smooth, with $n=3$ and
$\omega=\omega_{\mathrm{FS}}|_X$.
\[
 H^1(X;\C)=0\quad\leadsto\quad\text{Lefschetz hyperplane theorem},
 \qquad H^2(X;\C)=0.
\]
The Hodge diamond is
\[
\begin{array}{ccccccc}
 &&&1&&&\\
 &&0&&0&&\\
 &0&&1&&0&\\
 1&&101&&101&&1\\
 &0&&1&&0&\\
 &&0&&0&&\\
 &&&1&&&
\end{array}
\]
Hodge numbers stated without proof. The top and middle rows are
primitive; the other nonzero entries are $L$-images.

The complete string picture is
\[
 \underset{0}{1}\xrightarrow{L}
 \underset{2}{\omega}\xrightarrow{L}
 \underset{4}{\omega^2}\xrightarrow{L}
 \underset{6}{\omega^3},
\]
\[
 \Prim^3=\Hd^3,\qquad\text{dimension }204;\ \text{one-term strings}.
\]
\paragraph{Why are there no intermediate strings?}
\[
 \Hd^1=0,\qquad\Hd^2=\C\omega
 \quad\Longrightarrow\quad\Prim^1=\Prim^2=0.
\]
Every middle-degree harmonic form is primitive:
\[
 \Lambda:\Hd^3\longrightarrow\Hd^1=0,\qquad
 b_3=1+101+101+1=204.
\]
\textbf{One four-term string $+$ 204 one-term strings.}
All even cohomology is generated by $[\omega]$; all middle cohomology
is primitive.

% Slide 27.
\subsection*{Hard Lefschetz from primitive harmonic strings}
Fix $0\leq k\leq n$ and $0\leq r\leq\lfloor k/2\rfloor$.
Put $j=k-2r$, $m=n-k+r$.
\[
 L^{n-k}:L^r\Prim^j\xrightarrow{\sim}L^m\Prim^j,
 \qquad L^r\alpha\longmapsto L^m\alpha.
\]
\[
 \Hd^{2n-k}
 =\bigoplus_{r=0}^{\lfloor k/2\rfloor}L^{n-k+r}\Prim^{k-2r}.
\]
Thus
\[
 L^{n-k}:\Hd^k\xrightarrow{\sim}\Hd^{2n-k},\qquad0\leq k\leq n,
\]
\[
 L^r:\Hd^{n-r}\xrightarrow{\sim}\Hd^{n+r},\qquad0\leq r\leq n.
\]

% Slide 28.
\subsection*{Refinement by Hodge type}
Write $H^{p,q}(X)=H^{p,q}_{\dR}(X)$. Since $L$ has type $(1,1)$,
\[
 L:H^{p,q}(X)\longrightarrow H^{p+1,q+1}(X).
\]
For $k=p+q\leq n$, Hard Lefschetz therefore gives
\[
 L^{n-p-q}:H^{p,q}(X)\xrightarrow{\sim}H^{n-q,n-p}(X).
\]
\paragraph{Hodge decomposition $+$ Hard Lefschetz.}
Let $P^k(X)$ consist of classes with primitive harmonic
representatives, and set $P^{a,b}(X)=P^{a+b}(X)\cap H^{a,b}(X)$.
\[
 H^{p,q}(X)=\bigoplus_{r=\max(0,p+q-n)}^{\min(p,q)}
 L^rP^{p-r,q-r}(X).
\]

% Slide 29.
\subsection*{Review: the route to Hard Lefschetz}
\begin{enumerate}
\item K\"ahler identities: $[L,\Delta_d]=0$, $[\Lambda,\Delta_d]=0$.
Pointwise algebra: $\Lambda=L^*$, $[\Lambda,L]|_{A^k}=(n-k)\id$.
Analytic Hodge theorem: $\dim\Hd^k<\infty$,
$\phi_d^k:\Hd^k\xrightarrow{\sim}H^k_{\dR}$.
\item Lefschetz operators on harmonic forms: adjoint operators on the
finite-dimensional space $\Hd^{\bullet}$, with the degree commutator.
\item Primitive harmonic strings:
$\Prim^j=\ker\Lambda|_{\Hd^j}$,
$\alpha,L\alpha,\ldots,L^{n-j}\alpha$.
Exact lengths; injective coefficient maps.
\item Harmonic Lefschetz decomposition:
$\Hd^k=\bigoplus_rL^r\Prim^{k-2r}$.
Orthogonal, with unique coefficients.
\item Hard Lefschetz on harmonic forms:
$L^{n-k}:\Hd^k\xrightarrow{\sim}\Hd^{2n-k}$.
Matching primitive coefficients.
\item Hard Lefschetz on cohomology:
\[
 [\omega]^{n-k}\smile(-):H^k_{\mathrm{sing}}(X;\C)
 \xrightarrow{\sim}H^{2n-k}_{\mathrm{sing}}(X;\C).
\]
Transport by $\phi_d$ and de Rham.
\item Refinement by Hodge type:
$L^{n-p-q}:H^{p,q}\xrightarrow{\sim}H^{n-q,n-p}$.
Hodge decomposition yields the type refinement and primitive splitting.
\end{enumerate}

% Slide 30.
\section*{Consequences of Hard Lefschetz}
\paragraph{Growth towards the middle degree.}
For $k<n$, the map $L:H^k(X;\C)\longrightarrow H^{k+2}(X;\C)$
is injective:
\[
 L\alpha=0\quad\Longrightarrow\quad L^{n-k}\alpha=0
 \quad\Longrightarrow\quad\alpha=0.
\]
Since $L$ has bidegree $(1,1)$, this also holds on each Hodge type.
\[
 b_k\leq b_{k+2}\quad(k<n),\qquad
 h^{p,q}\leq h^{p+1,q+1}\quad(p+q<n).
\]
\paragraph{Every even degree has a nonzero K\"ahler-class power.}
For $k=0$, $H^0(X)\cong H^{2n}(X)$.
Hard Lefschetz gives $L^n(1)=[\omega]^n\neq0$.
If $[\omega]^r=0$, multiplying by $[\omega]^{n-r}$ would contradict this.
\[
 [\omega]^r\neq0,\qquad h^{r,r}\geq1,\qquad b_{2r}\geq1
 \quad(0\leq r\leq n).
\]

% Slide 31.
\section*{Hodge index for surfaces}
Let $X$ be a connected compact K\"ahler surface.
The intersection form
\[
 I:H^2(X;\R)\times H^2(X;\R)\longrightarrow\R,
 \qquad([\alpha],[\beta])\longmapsto\int_X\alpha\wedge\beta
\]
is symmetric. For $0\leq k\leq2$, Hard Lefschetz gives
\[
 H^k(X)=P^k(X)\oplus LH^{k-2}(X).
\]
For $n=k=2$,
\[
 H^2(X)=P^2(X)+LH^0(X),\qquad H^0(X)=\C,\qquad L[1]=[\omega].
\]
By the Lefschetz decomposition,
\[
 H^2(X)=\C[\omega]\orthsum P^2(X)\qquad(L^2).
\]
Refine by types:
\[
 H^{1,1}(X)_{\R}=\R[\omega]\orthsum P^{1,1}(X)_{\R}.
\]

% Slide 32. The handwritten final inequality in the norm identity is
% transcribed literally; no mathematical correction has been made.
\subsection*{Hodge index: orthogonality and signature}
\textbf{Goal.} The summands $\R[\omega]$ and $P^{1,1}(X)_{\R}$ are
$I$-orthogonal; $I$ is positive definite on the first and negative
definite on the second.

\paragraph{Orthogonality.}
For $\alpha\in\Prim^2$,
\[
 0=\langle\alpha,\omega\rangle_{L^2}
 =\int_X\alpha\wedge *\omega
 =\int_X\alpha\wedge\omega
 =I([\alpha],[\omega]),
\]
since $*\omega=\omega$. Thus
$\R[\omega]\perp_I P^{1,1}(X)_{\R}$.

\paragraph{The two signs.}
Let $\alpha\neq0$ be a primitive harmonic $(1,1)$-form.
\[
 *\alpha=-\alpha,
\]
\[
 \|\alpha\|_{L^2}^2=\int_X\alpha\wedge *\alpha
 =-I([\alpha],[\alpha])<0.
\]
\[
 dV_g=\omega^2/2,\qquad
 I([\omega],[\omega])=\int_X\omega^2=2\Vol_g(X)>0.
\]

% Slide 33.
\section*{Generalizing: Hodge--Riemann bilinear relations}
\paragraph{The Lefschetz bilinear form.}
Let $X$ be compact K\"ahler of complex dimension $n$, and fix
$0\leq k\leq n$. For $u=[\alpha]$ and $v=[\beta]$ in $H^k(X;\C)$,
define
\[
 Q_{\omega,k}(u,v)=(-1)^{k(k-1)/2}
 \int_X\alpha\wedge\beta\wedge\frac{\omega^{n-k}}{(n-k)!}.
\]
It depends only on the cohomology classes: change a representative by
an exact form and apply Stokes' theorem.
The pairing is $(-1)^k$-symmetric:
\[
 Q_{\omega,k}(v,u)=(-1)^kQ_{\omega,k}(u,v).
\]

% Slide 34.
\subsection*{Hodge--Riemann bilinear relations}
\textbf{Theorem.} On $P^k(X;\C)$, the form $Q_{\omega,k}$ is
nondegenerate, $(-1)^k$-symmetric, and pairs only conjugate Hodge types.
For $p+q=k$,
\[
 i^{p-q}Q_{\omega,k}(\alpha,\bar\alpha)>0
 \qquad\text{for }0\neq\alpha\in P^{p,q}(X).
\]
\paragraph{Proof sketch.}
For the primitive harmonic representative $\alpha$, positivity is the
norm identity:
\[
 \|\alpha\|_{L^2}^2
 =i^{p-q}Q_{\omega,k}([\alpha],[\bar\alpha]).
\]
For nondegeneracy, pair a nonzero Hodge component with its conjugate.
Other types contribute zero. Without primitivity, this positivity need
not hold.

\paragraph{Pointwise linear algebra.}
For a primitive $(p,q)$-form $\alpha$, with $k=p+q$,
\[
 *\bar\alpha=(-1)^{k(k-1)/2}i^{p-q}\bar\alpha
 \wedge\frac{\omega^{n-k}}{(n-k)!}.
\]
Substitute this into
$\|\alpha\|_{L^2}^2=\int_X\alpha\wedge *\bar\alpha$
to obtain the norm identity. Pointwise identity; its proof is omitted.

% Slide 35.
\section*{Polarized pure Hodge structures}
\paragraph{Pure of weight $k$.}
A finite-dimensional $\Q$-vector space $V_{\Q}$ with
\[
 V_{\C}:=V_{\Q}\otimes_{\Q}\C,\qquad
 V_{\C}=\bigoplus_{p+q=k}V^{p,q},\qquad
 \overline{V^{p,q}}=V^{q,p}.
\]
\paragraph{A rational polarization.}
A nondegenerate bilinear form $Q:V_{\Q}\times V_{\Q}\longrightarrow\Q$,
extended $\C$-bilinearly, such that
\[
 Q(\beta,\alpha)=(-1)^kQ(\alpha,\beta),
\]
\[
 Q(V^{p,q},V^{r,s})=0\quad\text{if }(r,s)\neq(q,p),
\]
\[
 i^{p-q}Q(\alpha,\bar\alpha)>0\qquad(0\neq\alpha\in V^{p,q}).
\]
For our compact K\"ahler $X$, assume $[\omega]\in H^2(X;\Q)$ and
$0\leq k\leq n$.
\[
 V_{\Q}=P^k(X;\Q)
 =\ker\bigl(L^{n-k+1}:H^k(X;\Q)\longrightarrow H^{2n-k+2}(X;\Q)\bigr).
\]
With its Hodge decomposition, this is polarized by the rational form
$Q_{\omega,k}$.

% Slide 36.
\section*{References and credits}
\begin{enumerate}
\item Claire Voisin. \emph{Hodge Theory and Complex Algebraic Geometry I.}
Cambridge Studies in Advanced Mathematics, vol.~76.
Cambridge University Press, 2002. Translated by Leila Schneps.
\item Raymond O. Wells, Jr. \emph{Differential Analysis on Complex Manifolds.}
Third edition. Graduate Texts in Mathematics, vol.~65.
Springer, 2008.
\item GPT 6. Assistance with slide preparation and LaTeX typesetting.
\end{enumerate}
\end{document}
